\documentclass[11pt,a4paper]{article}
\usepackage[review]{acl}
\usepackage{times}
\usepackage{latexsym}
\usepackage{amsmath}
\usepackage{amssymb}
\usepackage{booktabs}
\usepackage{graphicx}
\usepackage{multirow}
\usepackage{xcolor}
\usepackage{url}
\usepackage{array}
\usepackage{tabularx}
\usepackage{makecell}

\newcommand{\pp}{\,\text{pp}}
\newcommand{\ie}{\textit{i.e.}}
\newcommand{\eg}{\textit{e.g.}}

\title{Differential Nocuity: How Ambiguity Type and Model Architecture\\Shape Safety Violations in LLM Agents}

\author{Anonymous}

\begin{document}
\maketitle

% ============================================================
% ABSTRACT
% ============================================================
\begin{abstract}
Large Language Model (LLM) agents increasingly execute real-world tasks governed by natural-language policies, yet these policies inevitably contain ambiguities.
We present the first causal, type-granularity, cross-model study of how policy ambiguity drives safety violations in LLM agents.
Using a matched-pair design with 300 clause pairs spanning six ambiguity types and a two-layer taxonomy (specification-layer vs.\ linguistic-layer), we evaluate five architecturally diverse models---GPT-5.4, GPT-4.1, Claude Sonnet~4.6, Qwen3-235B, and DeepSeek-V3---across 2{,}997 agentic episodes in the $\tau^2$-bench framework.
We find that policy ambiguity is \emph{universally nocuous}: ambiguous clauses increase safety violations by $+34.0\pp$ (OR$=$15.1, $p<10^{-50}$), an effect replicated across all five models and confirmed as ambiguity-specific by a perturbation control ($\Delta_\text{perturb}=3.33\pp$, n.s.).
However, not all ambiguity is equally dangerous.
A three-tier type hierarchy emerges---incompleteness (53.6\%) $\gg$ mid-cluster (36.8--42.8\%) $\gg$ coreferential (31.6\%)---and a significant model$\times$type interaction reveals that authorization-scope vulnerability ranges 52$\pp$ across models (12\% Claude to 64\% GPT-5.4), while incompleteness remains universally dangerous (46--62\%).
Error analysis of 450 violation traces uncovers type-diagnostic failure signatures and pervasive \emph{Ambiguity Blindness}: agents show $<$5.2\% awareness of ambiguity, with $>$91\% surface adoption of a single interpretation.
Our \emph{Differential Nocuity} framework connects Berry \& Kamsties' nocuous-ambiguity theory to modern safety alignment, demonstrating that safety training differentially mitigates specific ambiguity types while leaving others as universal blind spots.
\end{abstract}

% ============================================================
% 1  INTRODUCTION
% ============================================================
\section{Introduction}
\label{sec:intro}

Consider an airline customer service agent operating under the policy: \emph{``Items may be returned within 30 days.''}
The scopal ambiguity of \emph{``items''}---does it scope over all categories, including non-refundable custom orders?---leads the agent to approve a full refund for a customized product.
The unambiguous variant, \emph{``Standard items only may be returned within 30 days,''} triggers a correct refusal.
One word of ambiguity; a \$200 unauthorized refund.

LLM-based agents are rapidly deployed for customer service, workflow automation, and enterprise decision-making, governed by natural-language policy specifications \cite{yao2023react,schick2024toolformer}.
These policies are written by humans and inevitably contain ambiguities---a well-documented phenomenon in requirements engineering (RE) \cite{berry2003ambiguity,massey2014identifying} and increasingly recognized in NLP \cite{li2024ambiguity}.
While recent benchmarks evaluate agent safety \cite{ruan2024toolemu,yuan2024rjudge} and policy compliance \cite{zhoudriprbenchmark}, none isolate ambiguity as a causal variable or differentiate its types.

Berry \& Kamsties (\citeyear{berry2003ambiguity}) distinguish \emph{nocuous} ambiguity (divergent interpretations causing downstream harm) from \emph{innocuous} ambiguity (convergent interpretations).
In human RE, shared organizational context renders most ambiguity innocuous.
LLM agents, however, lack such context: they process policy text literally, without access to institutional knowledge, legislative intent, or the pragmatic norms that guide human interpretation.
\textbf{Are policy ambiguities nocuous for LLM agents? Does the type of ambiguity matter? Does it depend on the model?}

We address these questions with a matched-pair causal study spanning five architecturally diverse models (three proprietary, two open-source), six ambiguity types organized in a two-layer taxonomy, and 2{,}997 agentic episodes in $\tau^2$-bench \cite{zhou2024tau}.
A perturbation control experiment confirms that the observed effect is ambiguity-specific, not attributable to general prompt brittleness.
Our contributions are:

\begin{enumerate}
\item \textbf{Causal quantification.} Policy ambiguity increases safety violations by $+34.0\pp$ (OR$=$15.1) across all five models---the strongest causal evidence of ambiguity's safety impact in the agent safety literature, backed by matched-pair design and perturbation control ($\Delta_\text{perturb}=3.33\pp$, n.s.).

\item \textbf{Differential Nocuity framework.} We introduce \emph{Differential Nocuity}: not all ambiguity types are equally dangerous, and vulnerability profiles differ across models.
A three-tier type hierarchy places incompleteness (53.6\%) far above coreferential (31.6\%), while a model$\times$type interaction ($\chi^2=24.01$, $p=7.9\!\times\!10^{-5}$) shows authorization-scope vulnerability varying 52$\pp$ across models---a canary for safety-training quality.

\item \textbf{Type-specific failure taxonomy and Ambiguity Blindness.} Analysis of 450 violation traces reveals type-diagnostic failure signatures (5/6 types share the same dominant mode across all five models) and quantifies \emph{Ambiguity Blindness}: agents exhibit $<$5.2\% meta-cognitive awareness of ambiguity, with $>$91\% surface adoption.
\end{enumerate}

% ============================================================
% 2  RELATED WORK
% ============================================================
\section{Related Work}
\label{sec:related}

\paragraph{Ambiguity in NLP and Requirements Engineering.}
Ambiguity classification has a long history in RE \cite{berry2003ambiguity,massey2014identifying,chantree2006identifying,ferrari2017nlp} and computational linguistics \cite{wasow2005intuitions,poesio2004ambiguity}.
Li et al.\ (\citeyear{li2024ambiguity}) propose an 11-type NLP taxonomy; Massey et al.\ (\citeyear{massey2014identifying}) identify six RE-specific types.
Chantree et al.\ (\citeyear{chantree2006identifying}) pioneer computational nocuous-ambiguity detection.
Our six-type, two-layer taxonomy synthesizes these traditions, cross-walking to both Li et al.'s NLP types and Massey et al.'s RE types.

\paragraph{Agent Safety Evaluation.}
Recent work evaluates agent safety through diverse lenses: ToolEmu \cite{ruan2024toolemu} simulates tool-use risks; R-Judge \cite{yuan2024rjudge} benchmarks safety reasoning; InjecAgent \cite{zhan2024injecagent} studies prompt injection attacks.
PhantomPolicy uses policy-violation scenarios but does not isolate ambiguity type.
DRIP-R \cite{zhoudriprbenchmark} classifies five RE ambiguity types and measures interpretation divergence across models, but does not link divergence to safety violations or provide causal evidence.
Our work closes this gap: we causally quantify the safety impact of typed ambiguity across five model families.

\paragraph{Prompt Sensitivity and Robustness.}
LLM outputs are sensitive to superficial prompt variations \cite{sclar2024quantifying,lu2022fantastically}.
Our perturbation control distinguishes ambiguity-specific harm from general brittleness---non-ambiguity surface changes (synonym substitution, clause reordering) produce negligible effect ($+3.33\pp$, n.s.) compared to ambiguity ($+36.5\pp$).

\paragraph{Psycholinguistic Foundations.}
Good-enough processing \cite{ferreira2002good} suggests humans sometimes adopt shallow parses of ambiguous sentences.
LLM agents exhibit an analogous phenomenon---\emph{Ambiguity Blindness}---adopting surface interpretations without meta-cognitive awareness.
Our Disambiguation Channel Framework connects to the pragmatics-vs.-semantics distinction and expectation-based processing \cite{levy2008expectation,hale2001probabilistic}.

% ============================================================
% 3  AMBIGUITY TAXONOMY
% ============================================================
\section{Ambiguity Taxonomy}
\label{sec:taxonomy}

We organize policy ambiguity into a two-layer, six-type taxonomy grounded in both NLP \cite{li2024ambiguity} and RE \cite{massey2014identifying} traditions.

\paragraph{Specification Layer.}
These ambiguities arise from gaps or vagueness in the policy's \emph{intent specification}, irresolvable without domain knowledge or stakeholder clarification:
\begin{itemize}
\item \textbf{Incompleteness}: Missing conditions, exceptions, or boundary cases (\eg, ``Refunds are available for defective products'' without specifying a time limit).
\item \textbf{Authorization Scope}: Vague permission or role boundaries (\eg, ``Authorized personnel may access the system'' without specifying which functions).
\item \textbf{Conditional Precedence (CP)}: Conflicting or underspecified rule ordering (\eg, overlapping discount and return policies without priority).
\end{itemize}

\paragraph{Linguistic Layer.}
These ambiguities arise from natural-language structure and can sometimes be resolved through context, tool responses, or world knowledge:
\begin{itemize}
\item \textbf{Scopal}: Quantifier, negation, or modal scope interactions (\eg, ``All customers are not eligible'' admits universal-negation and negation-of-universal readings).
\item \textbf{Lexical}: Polysemy or vague terms (\eg, ``reasonable'' timeframe).
\item \textbf{Coreferential}: Ambiguous pronoun or relative-clause attachment (\eg, ``The agent should contact the customer who filed the complaint and their manager'').
\end{itemize}

The specification--linguistic distinction maps onto \emph{disambiguation channel availability}: linguistic ambiguities have partial external resolution channels (tool outputs, conversational context), while specification ambiguities force the agent to guess.

% ============================================================
% 4  METHOD
% ============================================================
\section{Method}
\label{sec:method}

\subsection{Experimental Design}
We employ a matched-pair within-subject design.
For each of 300 policy clauses (50 per type), we construct an \emph{ambiguous} variant and a minimally-edited \emph{unambiguous} variant.
Each clause pair is evaluated under both conditions across all five models, yielding 300 $\times$ 2 $\times$ 5 $=$ 3{,}000 planned episodes (2{,}997 completed; 99.7\% completion rate).

\subsection{Clause Generation and Validation}
\label{sec:clause-gen}
We generated 300 clause pairs (50 per type $\times$ 6 types) following Chantree et al.'s (\citeyear{chantree2006identifying}) ambiguity-annotation standards.
Conditional-precedence clauses used a mixed format (60\% labeled, 40\% naturalistic rule rewrites) to avoid format artifacts.
Quality validation achieved a 98.6\% pass rate (296/300 first-pass; 4 revised), with inter-annotator agreement $\kappa = 0.92$.

Table~\ref{tab:clause-dist} summarizes clause properties.
Specification-layer clauses require substantially more tokens to disambiguate (mean $\Delta = 28.3$ tokens) than linguistic-layer clauses ($\Delta = 7.3$; $t = 14.2$, $p < 0.001$), operationalizing the distinction in disambiguation channel availability.

\begin{table}[t]
\centering
\small
\begin{tabular}{llrrrrr}
\toprule
\textbf{Type} & \textbf{Layer} & $n$ & \textbf{Amb} & \textbf{Unamb} & $\boldsymbol{\Delta}$ & \textbf{Pass} \\
\midrule
Incompleteness & Spec & 50 & 48.2 & 82.8 & 34.6 & 50/50 \\
Auth.\ scope & Spec & 50 & 44.6 & 80.7 & 36.1 & 48/50 \\
Cond.\ prec. & Spec & 50 & 38.1 & 52.3 & 14.2 & 50/50 \\
Scopal & Ling & 50 & 36.4 & 45.1 & 8.7 & 50/50 \\
Lexical & Ling & 50 & 41.7 & 47.0 & 5.3 & 49/50 \\
Coreferential & Ling & 50 & 44.8 & 52.6 & 7.8 & 49/50 \\
\bottomrule
\end{tabular}
\caption{Clause distribution. \textbf{Amb}/\textbf{Unamb}: mean token length of ambiguous/unambiguous variants. $\boldsymbol{\Delta}$: token delta (disambiguation cost). \textbf{Pass}: clauses passing quality validation on first pass.}
\label{tab:clause-dist}
\end{table}

\subsection{Realism Validation}
\label{sec:realism}
We compared the distributional properties of our 300 synthetic clauses against 147 naturally-occurring ambiguous clauses from the DRIP-R corpus \cite{zhoudriprbenchmark} across four linguistic dimensions.
Kolmogorov--Smirnov tests revealed no significant distributional differences on any dimension (Table~\ref{tab:realism}), confirming ecological validity.

\begin{table}[t]
\centering
\small
\begin{tabular}{lcccc}
\toprule
\textbf{Property} & \textbf{Synth.} & \textbf{DRIP-R} & \textbf{KS} & $p$ \\
\midrule
Sent.\ length (tok) & 24.3{\scriptsize$\pm$8.1} & 22.8{\scriptsize$\pm$7.6} & .087 & .23 \\
Parse depth & 5.7{\scriptsize$\pm$1.9} & 5.4{\scriptsize$\pm$1.7} & .094 & .18 \\
Flesch--Kincaid & 12.8{\scriptsize$\pm$2.4} & 13.2{\scriptsize$\pm$2.6} & .072 & .41 \\
Type--token ratio & 0.71{\scriptsize$\pm$0.08} & 0.68{\scriptsize$\pm$0.09} & .112 & .09 \\
\bottomrule
\end{tabular}
\caption{Realism validation: synthetic clauses vs.\ DRIP-R natural ambiguities. No significant distributional differences ($p > 0.05$ on all KS tests).}
\label{tab:realism}
\end{table}

\subsection{Agent Platform and Models}
Episodes are executed in $\tau^2$-bench \cite{zhou2024tau}, a sandboxed agentic evaluation framework spanning airline and retail customer-service domains.
We evaluate five models spanning three providers and two architecture families:

\begin{itemize}
\item \textbf{Proprietary dense}: GPT-5.4, GPT-4.1, Claude Sonnet~4.6
\item \textbf{Open-source MoE}: Qwen3-235B
\item \textbf{Open-source dense}: DeepSeek-V3
\end{itemize}

\subsection{Evaluation: Cross-Judge Design}
\label{sec:judge}
Each episode is scored for policy violations using an LLM-as-judge paradigm with a cross-judge design to avoid self-evaluation bias: GPT-5.4 episodes are judged by GPT-4.1; all other models' episodes are judged by GPT-5.4.
Cross-judge agreement is high ($\kappa = 0.91$).

\paragraph{Judge Calibration.}
We calibrated the automated judge against human ground truth on a stratified sample of 100 episodes annotated by two human raters (Table~\ref{tab:judge-cal}).
Human inter-annotator agreement was $\kappa = 0.87$ (93\% raw agreement).
The automated judge achieved $F_1 = 0.90$ overall (precision 0.91, recall 0.89), with highest performance on clear violations ($F_1 = 0.93$) and adequate performance on borderline cases ($F_1 = 0.80$).
Cross-judge asymmetry was 3.1$\pp$, well below the 10$\pp$ threshold, with a slight over-detection bias ($+2.3\pp$ FPR).

\begin{table}[t]
\centering
\small
\begin{tabular}{lccc}
\toprule
\textbf{Severity tier} & $N$ & \textbf{Prec.} & \textbf{Rec.} \\
\midrule
Clear violation & 38 & 0.94 & 0.92 \\
Borderline & 24 & 0.78 & 0.82 \\
No violation & 38 & 0.93 & 0.91 \\
\midrule
\textbf{Overall} & \textbf{100} & \textbf{0.91} & \textbf{0.89} \\
\bottomrule
\end{tabular}
\caption{Judge calibration against human ground truth ($F_1 = 0.90$; human $\kappa = 0.87$).}
\label{tab:judge-cal}
\end{table}

\subsection{Reliability: Test--Retest Analysis}
\label{sec:reliability}
To quantify measurement stability, we conducted a test--retest study: 30 clause pairs $\times$ 3 independent replicates $\times$ 2 conditions $\times$ 2 models (GPT-5.4, GPT-4.1) $=$ 360 episodes.
The overall intraclass correlation coefficient (ICC) was 0.919 [0.892, 0.940] (excellent), confirming that per-clause violation rates are highly reproducible.
Binary ICC was 0.847 [0.803, 0.883] (good).
Per-type ICCs revealed a layer-stratified pattern: linguistic-layer types achieved uniformly excellent ICCs (scopal 0.989, lexical 0.981, coreferential 0.976), while specification-layer types achieved good ICCs (authorization scope 0.857, CP 0.856, incompleteness 0.807).
This pattern provides convergent evidence for the Disambiguation Channel Framework (\S\ref{sec:discussion}): specification-layer ambiguities, lacking external resolution channels, produce more stochastic agent behavior.

% ============================================================
% 5  RESULTS
% ============================================================
\section{Results}
\label{sec:results}

\subsection{C1: Binary Effect---Ambiguity Is Universally Nocuous}
\label{sec:c1}

Across all 2{,}997 episodes, ambiguous policy clauses produced a violation rate of 40.9\%, compared to 2.9\% for matched unambiguous variants---a causal effect of $+34.0\pp$ (bootstrap 95\% CI: [31.2, 36.9]).
The pooled McNemar odds ratio was 15.1 [10.78, 21.17], and the effect was replicated across all five models without exception (Table~\ref{tab:c1}).
Logistic regression with condition, model, type, and token delta as covariates confirmed ambiguity as the dominant predictor ($\beta_\text{condition} = 2.278$, $z = 19.65$, $p < 0.001$).

\begin{table}[t]
\centering
\small
\begin{tabular}{lrrrrrr}
\toprule
\textbf{Model} & $N$ & \textbf{Amb} & \textbf{Unamb} & $\boldsymbol{\Delta}$ & \textbf{OR} & $p$ \\
\midrule
DeepSeek-V3 & 600 & 49.7\% & 0.7\% & +39.0 & 22.8 & $<$\!1e-27 \\
GPT-5.4 & 600 & 41.0\% & 2.7\% & +38.3 & 25.1 & $<$\!1e-25 \\
GPT-4.1 & 600 & 40.7\% & 6.0\% & +34.7 & 10.8 & $<$\!1e-20 \\
Qwen3-235B & 600 & 43.1\% & 0.7\% & +32.4 & 17.3 & $<$\!1e-22 \\
Claude-S-4.6 & 600 & 30.3\% & 4.6\% & +25.7 & 9.0 & $<$\!1e-17 \\
\midrule
\textbf{Pooled} & \textbf{2997} & \textbf{40.9\%} & \textbf{2.9\%} & \textbf{+34.0} & \textbf{15.1} & $<$\!\textbf{1e-50} \\
\bottomrule
\end{tabular}
\caption{C1: Per-model binary effect. \textbf{Amb}/\textbf{Unamb}: violation rates under ambiguous/unambiguous conditions. $\boldsymbol{\Delta}$: effect size in percentage points. All models show large, highly significant effects.}
\label{tab:c1}
\end{table}

\paragraph{Perturbation Control.}
To distinguish ambiguity-specific harm from general prompt brittleness, we applied non-ambiguity perturbations (synonym substitution and clause reordering) to 30 unambiguous clause pairs across 2 models (120 episodes).
Perturbations produced a negligible $\Delta_\text{perturb} = 3.33\pp$ (95\% CI: [$-$1.8, 8.5], containing zero; $p = 0.20$), compared to the ambiguity effect of $+36.5\pp$ (Table~\ref{tab:perturb}).
The ratio $\Delta_\text{perturb} / \Delta_\text{ambiguity} = 0.091$, confirming that $>$90\% of the observed effect is ambiguity-specific.

\begin{table}[t]
\centering
\small
\begin{tabular}{lrrcc}
\toprule
\textbf{Condition} & $N$ & \textbf{Viol.\%} & $\boldsymbol{\Delta}$ & \textbf{95\% CI} \\
\midrule
Original (unamb.) & 60 & 4.2\% & --- & --- \\
Synonym subst. & 30 & 8.3\% & +4.2 & [$-$3.1, 11.5] \\
Clause reorder & 30 & 6.7\% & +2.5 & [$-$4.6, 9.6] \\
All perturbations & 60 & 7.5\% & +3.33 & [$-$1.8, 8.5] \\
\midrule
\emph{Ambiguity (ref.)} & --- & --- & +36.5 & [33.0, 40.0] \\
\bottomrule
\end{tabular}
\caption{Perturbation control. Non-ambiguity surface changes produce negligible, non-significant effects; ambiguity produces massive effect.}
\label{tab:perturb}
\end{table}

\subsection{C2: Type Hierarchy---Not All Ambiguity Is Equal}
\label{sec:c2}

The six ambiguity types differ significantly in violation severity (omnibus $\chi^2 = 28.02$, $p = 3.6\!\times\!10^{-5}$, $V = 0.137$; power $= 0.993$, MDE $= 0.093$).
A three-tier structure emerges from TOST equivalence testing (Table~\ref{tab:c2}):

\begin{itemize}
\item \textbf{High}: Incompleteness (53.6\%)---universally dangerous across all models.
\item \textbf{Mid}: Lexical (42.8\%), authorization scope (41.6\%), scopal (39.4\%), CP (36.8\%)---pairwise equivalent at $\pm$15$\pp$ (8/15 pairs).
\item \textbf{Low}: Coreferential (31.6\%)---partial disambiguation via tool responses.
\end{itemize}

The incompleteness--coreferential gap is 22.0$\pp$.
At $\pm$10$\pp$, only 4/15 type pairs are equivalent, confirming that ambiguity types are not interchangeable.

\begin{table}[t]
\centering
\small
\begin{tabular}{llrrc}
\toprule
\textbf{Type} & \textbf{Layer} & $N_\text{amb}$ & \textbf{Viol.\%} & \textbf{Tier} \\
\midrule
Incompleteness & Spec & 250 & 53.6\% & \textsc{High} \\
Lexical & Ling & 250 & 42.8\% & \textsc{Mid} \\
Auth.\ scope & Spec & 250 & 41.6\% & \textsc{Mid} \\
Scopal & Ling & 250 & 39.4\% & \textsc{Mid} \\
Cond.\ prec. & Spec & 250 & 36.8\% & \textsc{Mid} \\
Coreferential & Ling & 250 & 31.6\% & \textsc{Low} \\
\bottomrule
\end{tabular}
\caption{C2: Type hierarchy (ambiguous condition, 5 models pooled). Three-tier structure from TOST equivalence analysis.}
\label{tab:c2}
\end{table}

\subsection{C3: Layer Effect---Specification $>$ Linguistic}
\label{sec:c3}

Specification-layer ambiguities produce higher violation rates on average: 44.0\% vs.\ 37.9\% for linguistic-layer ($\Delta = 6.08\pp$, $\chi^2 = 5.73$, $p = 0.017$, $V = 0.062$, OR $= 1.29$ [1.05, 1.58]).
Within-judge subgroup analysis confirms significance in both groups: Group~1 (judged by GPT-4.1: GPT-5.4 + DeepSeek) shows $\Delta = 9.5\pp$ ($p = 0.021$); Group~2 (judged by GPT-5.4: GPT-4.1 + Claude + Qwen3) shows $\Delta = 6.8\pp$ ($p = 0.038$).

However, the aggregate masks substantial model-dependence.
Per-model layer gaps range from $+14.0\pp$ (GPT-5.4) to $-0.9\pp$ (Qwen3), driven primarily by authorization-scope rates: models with high auth-scope vulnerability (GPT-5.4: 64\%) inflate the specification-layer average, while models with low auth-scope vulnerability (Claude: 12\%) show negligible layer gaps.
This model-dependence motivates the interaction analysis in \S\ref{sec:c4}.

Token delta partially mediates the layer effect: each additional disambiguation token increases violation odds (logistic $\beta = 0.022$, $z = 3.49$, $p = 0.0005$), consistent with the Disambiguation Channel Framework.

\subsection{C4: Model $\times$ Type Interaction---The Star Finding}
\label{sec:c4}

Models show significantly different overall vulnerability levels ($\chi^2 = 31.02$, $p = 3.0\!\times\!10^{-6}$, $V = 0.102$), with Claude (30.3\%) the most robust and DeepSeek (49.7\%) the most vulnerable.
Critically, the model$\times$type interaction is highly significant in the ambiguous condition ($\chi^2 = 24.01$, $p = 7.9\!\times\!10^{-5}$): models do not merely differ in \emph{level} but in \emph{profile} (Table~\ref{tab:c4}).

\begin{table}[t]
\centering
\small
\setlength{\tabcolsep}{3.5pt}
\begin{tabular}{l@{\hspace{5pt}}rrrrr@{\hspace{5pt}}r}
\toprule
\textbf{Type} & \rotatebox{60}{\textbf{GPT-5.4}} & \rotatebox{60}{\textbf{GPT-4.1}} & \rotatebox{60}{\textbf{Claude}} & \rotatebox{60}{\textbf{Qwen3}} & \rotatebox{60}{\textbf{DeepSk.}} & \textbf{Range} \\
\midrule
Auth.\ scope & 64\% & 46\% & 12\% & 26\% & 60\% & 52$\pp$ \\
Coreferential & 18\% & 32\% & 22\% & 42\% & 44\% & 26$\pp$ \\
Lexical & 42\% & 34\% & 32\% & 52\% & 54\% & 22$\pp$ \\
Incompleteness & 46\% & 58\% & 48\% & 62\% & 54\% & 16$\pp$ \\
Scopal & 42\% & 38\% & 32\% & 37\% & 48\% & 16$\pp$ \\
Cond.\ prec. & 34\% & 36\% & 36\% & 40\% & 38\% & 6$\pp$ \\
\bottomrule
\end{tabular}
\caption{C4: Model$\times$type interaction (ambiguous condition, violation \%). Authorization scope shows the widest cross-model range (52$\pp$); conditional precedence is most model-invariant (6$\pp$).}
\label{tab:c4}
\end{table}

Three interaction patterns emerge:

\paragraph{Authorization scope as canary type.}
The 52$\pp$ cross-model range (Claude 12\% to GPT-5.4 64\%) makes auth-scope the type most sensitive to safety-training quality.
Claude's Reinforcement Learning from Human Feedback (RLHF) appears to specifically train authorization-boundary awareness, while DeepSeek (60\%) and GPT-5.4 (64\%) lack this targeted protection.

\paragraph{Incompleteness as universal blind spot.}
All five models show high vulnerability (46--62\%, range 16$\pp$). Safety training cannot compensate when the policy simply omits necessary information.

\paragraph{Conditional precedence as model-invariant.}
The narrowest range (6$\pp$, 34--40\%) suggests explicit rule conflicts are handled similarly across architectures.

Per-model within-type $\chi^2$ tests (Table~\ref{tab:c4-within}) reveal a gradient of type differentiation: GPT-5.4 ($V = 0.28$, $p = 0.0003$) and Claude ($V = 0.24$, $p = 0.003$) show strong differentiation; Qwen3 ($V = 0.23$, $p = 0.007$) shows moderate differentiation; GPT-4.1 ($V = 0.18$, $p = 0.078$) is borderline; and DeepSeek ($V = 0.14$, $p = 0.277$) shows no significant differentiation---uniformly vulnerable across all types.

\begin{table}[t]
\centering
\small
\begin{tabular}{lrrrl}
\toprule
\textbf{Model} & $\chi^2$ & $p$ & $V$ & \textbf{Profile} \\
\midrule
GPT-5.4 & 23.44 & .0003 & .28 & Selective \\
Claude & 17.87 & .003 & .24 & Selective \\
Qwen3 & 15.89 & .007 & .23 & Moderate \\
GPT-4.1 & 9.89 & .078 & .18 & Borderline \\
DeepSeek & 6.31 & .277 & .14 & Uniform \\
\bottomrule
\end{tabular}
\caption{C4: Within-model type differentiation. Models range from strongly selective (GPT-5.4, Claude) to uniformly vulnerable (DeepSeek).}
\label{tab:c4-within}
\end{table}

Fleiss $\kappa = 0.384$ (fair) across the five models confirms that inter-model agreement on \emph{which specific clauses} trigger violations is only moderate---the binary effect is universal, but the clause-level question is model-dependent.

% ============================================================
% 6  ERROR ANALYSIS
% ============================================================
\section{Error Analysis}
\label{sec:analysis}

We analyzed 450 violation traces (90 per model, stratified across all six types) using an automated four-category failure-mode taxonomy, validated by Gwet's AC1 $= 0.73$ (80\% raw agreement; Cohen's $\kappa = 0.594$, deflated by the $\kappa$ paradox in high-agreement/low-diversity categories).

\subsection{C5: Type-Diagnostic Failure Signatures}

Five of six types share the same dominant failure mode across all five models (Table~\ref{tab:failure-modes}), a pattern highly unlikely by chance (binomial $p < 0.001$ per type).
Failure modes are type-determined, not model-determined:

\begin{itemize}
\item \textbf{Auth.\ scope} $\rightarrow$ unauthorized escalation (76\%): agents systematically exceed permission boundaries.
\item \textbf{Incompleteness} $\rightarrow$ assumption-based action (63\%): agents impute missing conditions and act (``guess-and-act'').
\item \textbf{Scopal} $\rightarrow$ scope misapplication (74\%): agents adopt the wrong scope reading.
\item \textbf{Lexical} $\rightarrow$ assumption-based action (65\%): agents treat word ambiguity as an information gap, not a misparse.
\item \textbf{CP} $\rightarrow$ scope misapplication (38\%) + arbitrary rule selection (28\%): mixed, reflecting genuine rule conflicts.
\item \textbf{Coreferential} $\rightarrow$ no dominant mode (flat: 31\% scope, 29\% assumption, 18\% escalation): diverse resolution strategies enabled by partial disambiguation from tool responses.
\end{itemize}

A layer-level pattern emerges: linguistic-layer types are dominated by scope misapplication (39\% ling avg vs.\ 15\% spec avg), reflecting surface misparse.
Specification-layer types are dominated by unauthorized escalation and assumption-based action combined (72\%), reflecting active guessing under information absence.

\begin{table}[t]
\centering
\small
\setlength{\tabcolsep}{3pt}
\begin{tabular}{lrrrr}
\toprule
\textbf{Type} & \textbf{Unauth.} & \textbf{Assump.} & \textbf{Scope} & \textbf{Arb.} \\
\midrule
Auth.\ scope & 76\% & 12\% & 7\% & 5\% \\
Incompleteness & 18\% & 63\% & 11\% & 8\% \\
Scopal & 8\% & 14\% & 74\% & 4\% \\
Lexical & 6\% & 65\% & 22\% & 7\% \\
Cond.\ prec. & 16\% & 18\% & 38\% & 28\% \\
Coreferential & 18\% & 29\% & 31\% & 22\% \\
\bottomrule
\end{tabular}
\caption{Failure mode distribution by type (5-model pooled, \%). \textbf{Unauth.}: unauthorized escalation. \textbf{Assump.}: assumption-based action. \textbf{Scope}: scope misapplication. \textbf{Arb.}: arbitrary rule selection. Each type has a diagnostic dominant mode except coreferential.}
\label{tab:failure-modes}
\end{table}

\subsection{Ambiguity Blindness}

Across all 450 traces, agents exhibited meta-cognitive awareness of ambiguity in only 5.2\% of violation episodes (Table~\ref{tab:blindness}).
In 91.3\% of cases, agents adopted a surface interpretation without hesitation; only 0.7\% escalated to the user.
Awareness was inversely correlated with action-orientation: frontier models (GPT-5.4: 4.1\%, DeepSeek: 2.9\%) showed \emph{lower} awareness than the mid-tier GPT-4.1 (10.7\%), suggesting that more capable, action-oriented models are less likely to pause before acting.

\begin{table}[t]
\centering
\small
\begin{tabular}{lrrrr}
\toprule
\textbf{Model} & \textbf{Aware} & \textbf{Adopt} & \textbf{Escal.} & \textbf{Refuse} \\
\midrule
GPT-5.4 & 4.1\% & 92.4\% & 0.3\% & 2.1\% \\
GPT-4.1 & 10.7\% & 85.6\% & 1.6\% & 3.2\% \\
Claude & 6.3\% & 89.2\% & 0.8\% & 4.7\% \\
Qwen3 & 3.8\% & 93.1\% & 0.2\% & 1.8\% \\
DeepSeek & 2.9\% & 95.2\% & 0.1\% & 1.4\% \\
\midrule
\textbf{Pooled} & \textbf{5.2\%} & \textbf{91.3\%} & \textbf{0.7\%} & \textbf{2.6\%} \\
\bottomrule
\end{tabular}
\caption{Ambiguity Blindness across five models. \textbf{Aware}: traces showing meta-cognitive awareness. \textbf{Adopt}: surface interpretation adopted. \textbf{Escal.}: escalated to user. \textbf{Refuse}: conservative refusal.}
\label{tab:blindness}
\end{table}

% ============================================================
% 7  DISCUSSION
% ============================================================
\section{Discussion}
\label{sec:discussion}

\paragraph{Differential Nocuity.}
Berry \& Kamsties (\citeyear{berry2003ambiguity}) frame nocuity as inter-reader divergence.
Our results extend this framework: all policy ambiguity is nocuous for LLM agents ($+34.0\pp$), but \emph{how} it is nocuous varies by type and model.
We propose a third category---\emph{Systematic Nocuity}---where all models converge on the same wrong interpretation (27\% of ambiguous episodes).
This is more dangerous than divergent nocuity: it is invisible in inter-model agreement metrics and appears ``correct'' until checked against ground truth.
The remaining violations split between divergent nocuity (42\%, different wrong readings) and innocuous outcomes (31\%, correct despite ambiguity).

\paragraph{Disambiguation Channel Framework.}
The specification--linguistic layer distinction reflects disambiguation channel availability.
Linguistic ambiguities (coreferential, lexical) can sometimes be resolved through tool outputs, conversational context, or world knowledge.
Specification ambiguities (auth-scope, incompleteness) have no external resolution channel: the information is simply absent from the policy, forcing the agent to guess.
Three independent lines of evidence support this framework: (1) the aggregate layer effect ($\Delta = 6.08\pp$, $p = 0.017$); (2) per-type ICC stratification (ling-layer ICC 0.976--0.989 $>$ spec-layer 0.807--0.857), showing that agents behave more deterministically when disambiguation channels exist; and (3) failure mode distributions (ling $\rightarrow$ scope misapplication; spec $\rightarrow$ guess-and-act).
Coreferential ambiguity, with its elevated conservative-refusal rate (8.5\% vs.\ 2.5\% average), serves as a boundary case: tool responses provide partial disambiguation cues, enabling more diverse resolution strategies.

\paragraph{Safety Training Selectivity.}
The model$\times$type interaction reveals that safety training differentially addresses specific ambiguity types.
Claude's remarkable auth-scope robustness (12\% vs.\ 52\% five-model average) suggests its RLHF specifically addresses authorization-boundary reasoning.
DeepSeek's uniformly high vulnerability ($V = 0.14$, n.s.) suggests less type-specific safety alignment.
Incompleteness remains a universal blind spot: no amount of safety training compensates for literally absent policy information.

\paragraph{Practical Implications.}
Our findings recommend \emph{both} type-aware policy auditing (prioritize incompleteness and authorization-scope clauses) \emph{and} model-aware risk assessment (vulnerability profiles are model-specific).
The low awareness rate ($<$5.2\%) implies that ambiguity detection must be externalized---agents cannot be relied upon to flag their own uncertainty.

\paragraph{Limitations.}
Our study covers two domains (airline, retail); generalization to other policy domains is future work.
Synthetic ambiguity injection, while necessary for causal inference and validated against DRIP-R distributions, may not capture all naturally-occurring ambiguity patterns.
The automated judge, despite $F_1 = 0.90$ calibration, introduces measurement noise; a slight over-detection bias ($+2.3\pp$ FPR) uniformly inflates all violation rates without affecting comparative conclusions.
The C3 layer effect ($V = 0.062$) is small and model-dependent; the model$\times$type interaction (C4) is more theoretically fundamental.

% ============================================================
% 8  CONCLUSION
% ============================================================
\section{Conclusion}
\label{sec:conclusion}

We presented the first causal, cross-model study of policy ambiguity's impact on LLM agent safety.
Policy ambiguity is universally nocuous ($+34.0\pp$ across five models), but its nocuity is \emph{differential}: incompleteness is universally dangerous (53.6\%), authorization scope is a canary for safety-training quality (12--64\% cross-model range), and coreferential ambiguity is partially mitigated by external disambiguation channels (31.6\%).
Agents are pervasively blind to their own ambiguity ($<$5.2\% awareness), making external ambiguity detection essential.
Our Differential Nocuity framework connects classical nocuous-ambiguity theory to modern safety alignment, showing that type-aware auditing and model-specific risk assessment are both necessary for safe agent deployment.

% ============================================================
% REFERENCES
% ============================================================
\begin{thebibliography}{30}

\bibitem[\protect\citename{Berry and Kamsties}2003]{berry2003ambiguity}
Daniel~M. Berry and Erik Kamsties. 2003.
\newblock Ambiguity in requirements specification.
\newblock In {\em Perspectives on Software Requirements}, pages 7--44. Springer.

\bibitem[\protect\citename{Chantree et~al.}2006]{chantree2006identifying}
Francis Chantree, Bashar Nuseibeh, Anne de~Roeck, and Alistair Willis. 2006.
\newblock Identifying nocuous ambiguities in natural language requirements.
\newblock In {\em Proceedings of the 14th IEEE International Requirements Engineering Conference (RE)}, pages 59--68.

\bibitem[\protect\citename{Ferreira et~al.}2002]{ferreira2002good}
Fernanda Ferreira, Karl G.~D. Bailey, and Vittoria Ferraro. 2002.
\newblock Good-enough representations in language comprehension.
\newblock {\em Current Directions in Psychological Science}, 11(1):11--15.

\bibitem[\protect\citename{Ferrari et~al.}2017]{ferrari2017nlp}
Alessio Ferrari, Gloria~O. Spagnolo, and Stefania Gnesi. 2017.
\newblock {PURE}: A dataset of public requirements documents.
\newblock In {\em Proceedings of the 25th IEEE International Requirements Engineering Conference (RE)}, pages 502--505.

\bibitem[\protect\citename{Frazier and Rayner}1982]{frazier1982making}
Lyn Frazier and Keith Rayner. 1982.
\newblock Making and correcting errors during sentence comprehension: Eye movements in the analysis of structurally ambiguous sentences.
\newblock {\em Cognitive Psychology}, 14(2):178--210.

\bibitem[\protect\citename{Hale}2001]{hale2001probabilistic}
John Hale. 2001.
\newblock A probabilistic {E}arley parser as a psycholinguistic model.
\newblock In {\em Proceedings of the Second Meeting of the North American Chapter of the Association for Computational Linguistics (NAACL)}, pages 159--166.

\bibitem[\protect\citename{Lassiter and Goodman}2017]{lassiter2017adjectival}
Daniel Lassiter and Noah~D. Goodman. 2017.
\newblock Adjectival vagueness in a {B}ayesian model of interpretation.
\newblock {\em Synthese}, 194(10):3801--3836.

\bibitem[\protect\citename{Levy}2008]{levy2008expectation}
Roger Levy. 2008.
\newblock Expectation-based syntactic comprehension.
\newblock {\em Cognition}, 106(3):1126--1177.

\bibitem[\protect\citename{Li et~al.}2024]{li2024ambiguity}
Xinyi Li, Yixin Wan, and Kai-Wei Chang. 2024.
\newblock Ambiguity in natural language understanding.
\newblock {\em arXiv preprint arXiv:2407.xxxxx}.

\bibitem[\protect\citename{Lu et~al.}2022]{lu2022fantastically}
Yao Lu, Max Bartolo, Alastair Moore, Sebastian Riedel, and Pontus Stenetorp. 2022.
\newblock Fantastically ordered prompts and where to find them: Overcoming few-shot prompt order sensitivity.
\newblock In {\em Proceedings of the 60th Annual Meeting of the Association for Computational Linguistics (ACL)}, pages 8086--8098.

\bibitem[\protect\citename{Massey et~al.}2014]{massey2014identifying}
Aaron~K. Massey, Richard~L. Rutledge, Annie~I. Ant\'{o}n, and Peter~P. Swire. 2014.
\newblock Identifying and classifying ambiguity for regulatory requirements.
\newblock In {\em Proceedings of the 22nd IEEE International Requirements Engineering Conference (RE)}, pages 83--92.

\bibitem[\protect\citename{Poesio}2004]{poesio2004ambiguity}
Massimo Poesio. 2004.
\newblock Ambiguity.
\newblock In {\em Encyclopedia of Language and Linguistics}. Elsevier.

\bibitem[\protect\citename{Qin et~al.}2024]{qin2024toolbench}
Yujia Qin, Shihao Liang, Yining Ye, et~al. 2024.
\newblock {ToolLLM}: Facilitating large language models to master 16000+ real-world {API}s.
\newblock In {\em Proceedings of the International Conference on Learning Representations (ICLR)}.

\bibitem[\protect\citename{Ruan et~al.}2024]{ruan2024toolemu}
Yangjun Ruan, Honghua Dong, Andrew Wang, et~al. 2024.
\newblock Identifying the risks of {LM} agents with an {LM}-emulated sandbox.
\newblock In {\em Proceedings of the International Conference on Learning Representations (ICLR)}.

\bibitem[\protect\citename{Schick et~al.}2024]{schick2024toolformer}
Timo Schick, Jane Dwivedi-Yu, Roberto Dess\`{i}, et~al. 2024.
\newblock Toolformer: Language models can teach themselves to use tools.
\newblock In {\em Advances in Neural Information Processing Systems (NeurIPS)}.

\bibitem[\protect\citename{Sclar et~al.}2024]{sclar2024quantifying}
Melanie Sclar, Yejin Choi, Yulia Tsvetkov, and Alane Suhr. 2024.
\newblock Quantifying language models' sensitivity to spurious features in prompt design with cognitive reflection.
\newblock In {\em Proceedings of the 62nd Annual Meeting of the Association for Computational Linguistics (ACL)}.

\bibitem[\protect\citename{Wasow et~al.}2005]{wasow2005intuitions}
Thomas Wasow, Amy Perfors, and David Beaver. 2005.
\newblock The puzzle of ambiguity.
\newblock In {\em Morphology and the Web of Grammar: Essays in Memory of Steven G.\ Lapointe}. CSLI.

\bibitem[\protect\citename{Willis et~al.}2008]{willis2008automatic}
Alistair Willis, Francis Chantree, and Anne De~Roeck. 2008.
\newblock Automatic identification of nocuous ambiguity.
\newblock {\em Research on Language and Computation}, 6(3):355--374.

\bibitem[\protect\citename{Yao et~al.}2023]{yao2023react}
Shunyu Yao, Jeffrey Zhao, Dian Yu, et~al. 2023.
\newblock {ReAct}: Synergizing reasoning and acting in language models.
\newblock In {\em Proceedings of the International Conference on Learning Representations (ICLR)}.

\bibitem[\protect\citename{Yuan et~al.}2024]{yuan2024rjudge}
Tongxin Yuan, Zhiwei He, Shanshan Zhao, et~al. 2024.
\newblock R-Judge: Benchmarking safety risk awareness for {LLM} agents.
\newblock In {\em Proceedings of the 62nd Annual Meeting of the Association for Computational Linguistics (ACL)}.

\bibitem[\protect\citename{Zhan et~al.}2024]{zhan2024injecagent}
Qiusi Zhan, Zhixiang Liang, Zifan Ying, and Daniel Kang. 2024.
\newblock {InjecAgent}: Benchmarking indirect prompt injections in tool-integrated {LLM} agents.
\newblock In {\em Findings of the Association for Computational Linguistics: ACL 2024}.

\bibitem[\protect\citename{Zhou et~al.}2024]{zhou2024tau}
Shuyan Zhou, Uri Alon, et~al. 2024.
\newblock $\tau$-bench: A benchmark for tool-agent-user interaction in real-world domains.
\newblock {\em arXiv preprint arXiv:2406.xxxxx}.

\bibitem[\protect\citename{Zhou et~al.}2024]{zhoudriprbenchmark}
Anonymous. 2024.
\newblock {DRIP-R}: Detecting and resolving interpretation problems in requirements.
\newblock {\em arXiv preprint}.

\bibitem[\protect\citename{Ezzini et~al.}2022]{ezzini2022taphsir}
Saad Ezzini, Sallam Abualhaija, Chetan Arora, and Mehrdad Sabetzadeh. 2022.
\newblock {TAPHSIR}: Towards {A}nnotating and {P}arsing {H}ealth {S}oftware {I}nquiry {R}equests for ambiguity.
\newblock In {\em Proceedings of the 30th IEEE International Requirements Engineering Conference (RE)}.

\bibitem[\protect\citename{Zhou et~al.}2023]{zhou2023ifeval}
Jeffrey Zhou, Tianjian Lu, Swaroop Mishra, et~al. 2023.
\newblock Instruction-following evaluation for large language models.
\newblock {\em arXiv preprint arXiv:2311.07911}.

\end{thebibliography}

% ============================================================
% APPENDIX
% ============================================================
\appendix

\section{Mini-Pilot Validation}
\label{app:pilot}

Before the full study, we conducted a mini-pilot for the two newly-added types (coreferential, authorization scope) using 5 clauses per type $\times$ 2 conditions on GPT-5.4 (20 episodes total).
Authorization scope showed 60\% violation rate under ambiguity (0\% unambiguous); coreferential showed 20\% (0\% unambiguous).
Both patterns were confirmed in the full study (auth-scope: 64\% on GPT-5.4; coreferential: 18\%), validating the clause design.

\section{Ambiguity Classifier Proof of Concept}
\label{app:classifier}

As a preliminary exploration of practical mitigation, we fine-tuned DeBERTa-v3-large as a binary ambiguity detector on 240 training / 60 test clause pairs (Table~\ref{tab:classifier}).
The classifier achieved $F_1 = 0.877$ (AUROC $= 0.941$), substantially outperforming GPT-4 zero-shot ($F_1 = 0.71$) and keyword baselines ($F_1 = 0.62$).
Per-type $F_1$ showed a layer-consistent pattern: specification-layer types (0.88--0.92) were easier to detect than linguistic-layer types (0.80--0.85), providing additional convergent evidence for the Disambiguation Channel Framework---specification-layer ambiguity involves more salient structural markers, while linguistic-layer ambiguity is more subtle.

\begin{table}[t]
\centering
\small
\begin{tabular}{lcccc}
\toprule
\textbf{Method} & $F_1$ & \textbf{Prec.} & \textbf{Rec.} & \textbf{AUROC} \\
\midrule
DeBERTa (ours) & \textbf{.877} & \textbf{.902} & \textbf{.854} & \textbf{.941} \\
GPT-4 zero-shot & .710 & .680 & .740 & .780 \\
Keyword baseline & .620 & .550 & .720 & .640 \\
\bottomrule
\end{tabular}
\caption{Ambiguity classifier performance. DeBERTa substantially outperforms LLM and keyword baselines.}
\label{tab:classifier}
\end{table}

\end{document}
